Eight systems are compared on the Tennessee Eastman Process under one
protocol: a trivial floor, three classical models, two networks, and the
copilot with its ablation, all scored on validation folds of a partition
frozen by simulation run, over 15 folds from three seeds, and the same 15
models once more on the held-out test set. The proposed copilot reaches $0.741
\pm 0.008$ macro-averaged $F_1$ and $0.845 \pm 0.009$ Recall@3 on validation,
and $0.719$ macro-$F_1$ on test. Against the best classical model, logistic
regression at $0.652 \pm 0.005$, the paired difference is $+0.089 \pm 0.004$
over the seeds, which is larger than either dispersion and keeps its sign in
all 15 folds, and it is $+0.105 \pm 0.004$ on test, again in all 15. The
ablation separates the two things the agent does: the loop is worth more than
$0.05$ macro-$F_1$ over the same network scored once on validation and more
than $0.04$ on test, with the same sign in all 15 folds of both, while
retrieval costs about $0.01$ on validation, so the proposed method does not
beat its own ablation there, and is ahead of it on test by $+0.005 \pm 0.004$
over the seeds, inside the dispersion. It stores 19,601 numbers, trains
nothing because it loads the network saved for each fold, answers in under a
millisecond per episode, and calls no language model, so its token count is
zero. The main limitation is coverage rather than accuracy: requiring the
classifier and the retrieval to agree makes the copilot hand 59\% of episodes
to the operator against 26\% for the ablation, while accuracy on the episodes
it does answer goes from 96.0\% to 97.1\%; on test it hands over 59\% against
18\% and goes from 85.4\% to 95.1\%.
